% =============================================================================
% PREÂMBULO COMUM — Compartilhado entre light.tex e dark.tex
% =============================================================================
% Codificação e Idioma
\usepackage[utf8]{inputenc}
\usepackage[T1]{fontenc}
\usepackage{helvet}
\renewcommand{\familydefault}{\sfdefault}
\usepackage{csquotes}

% Índice Remissivo
\usepackage{makeidx}
\makeindex

% Controle de Overfull Boxes e Hifenização
\usepackage{microtype}
\tolerance=2000
\emergencystretch=3em
\hfuzz=2pt

% Numeração de Seções e Sumário até Nível 5
\setcounter{secnumdepth}{5}
\setcounter{tocdepth}{5}

% Padronização de Tabelas
\usepackage{etoolbox}
\AtBeginEnvironment{tabular}{\small}
\AtBeginEnvironment{longtable}{\small}
\AtBeginEnvironment{tabularx}{\small}
\hfuzz=2pt
\widowpenalty=10000
\clubpenalty=10000
\raggedbottom
\usepackage{seqsplit}
\usepackage{adjustbox}

\let\origtexttt\texttt

\newcommand{\codigoquebravel}[1]{\seqsplit{\origtexttt{#1}}}

% Matemática
\usepackage{amsmath,amssymb,amsfonts,amsthm}
\usepackage{mathtools}
\usepackage{bm}
\usepackage{mathrsfs}
\IfFileExists{dsfont.sty}{\usepackage{dsfont}}{}
\providecommand{\mathds}[1]{\mathbb{#1}}
\IfFileExists{stmaryrd.sty}{\usepackage{stmaryrd}}{}

% Gráficos e Ilustrações
\usepackage{graphicx}
\usepackage{tikz}
\usepackage{tikz-3dplot}
\usetikzlibrary{positioning,shapes,arrows,fit,backgrounds,calc}
\usetikzlibrary{decorations.pathmorphing,patterns,shadows}
\usetikzlibrary{automata,petri,topaths}
\usetikzlibrary{chains,scopes}
\usepackage{pgfplots}
\pgfplotsset{compat=1.18}
\usepackage{float}
\usepackage{caption}
\usepackage{subcaption}
\usepackage{wrapfig}

% Tabelas
\usepackage{booktabs}
\usepackage{longtable}
\usepackage{tabularx}
\usepackage{multirow}
\usepackage{array}
\usepackage{colortbl}
\usepackage{enumitem}

% Cores
\usepackage[table]{xcolor}

% Hiperlinks
\usepackage[alf]{abntex2cite}
\usepackage{hyperref}
\usepackage{url}

% Tradução de autorefs para português
\renewcommand{\chapterautorefname}{Capítulo}
\renewcommand{\tableautorefname}{Tabela}
\renewcommand{\partautorefname}{Parte}
\renewcommand{\sectionautorefname}{Seção}
\renewcommand{\subsectionautorefname}{Subseção}
\renewcommand{\subsubsectionautorefname}{Subsubseção}
\renewcommand{\paragraphautorefname}{Parágrafo}
\renewcommand{\figureautorefname}{Figura}

% Quebra de URLs
\def\UrlBreaks{\do\:\do\-\do\_\do\/\do\.\do\~\do\?\do\#\do\&\do\%\do\@\do\\}

% Código Fonte
\usepackage{listings}
\providecommand{\lstlistingautorefname}{Código}
\usepackage{verbatim}

\lstdefinestyle{pythonstyle}{
    language=Python,
    backgroundcolor=\color{codeback},
    commentstyle=\color{codegreen},
    keywordstyle=\color{codeblue},
    numberstyle=\tiny\color{codegray},
    stringstyle=\color{codepurple},
    basicstyle=\ttfamily\scriptsize,
    breakatwhitespace=false,
    breaklines=true,
    postbreak=\mbox{$\hookrightarrow$\space},
    captionpos=b,
    keepspaces=true,
    numbers=left,
    numbersep=5pt,
    showspaces=false,
    showstringspaces=false,
    showtabs=false,
    tabsize=4,
    frame=single,
    rulecolor=\color{coderule},
    frameround=tttt,
    emph={self,True,False,None},
    emphstyle=\color{magenta},
    extendedchars=true,
    literate={á}{{\'a}}1 {é}{{\'e}}1 {í}{{\'i}}1 {ó}{{\'o}}1 {ú}{{\'u}}1
             {Á}{{\'A}}1 {É}{{\'E}}1 {Í}{{\'I}}1 {Ó}{{\'O}}1 {Ú}{{\'U}}1
             {ã}{{\~a}}1 {õ}{{\~o}}1 {Ã}{{\~A}}1 {Õ}{{\~O}}1
             {ç}{{\c{c}}}1 {Ç}{{\c{C}}}1 {â}{{\^a}}1 {ê}{{\^e}}1
             {×}{{\ensuremath{\times}}}1 {—}{{\textemdash}}1
}

\lstdefinestyle{bashstyle}{
    language=bash,
    backgroundcolor=\color{codeback},
    commentstyle=\color{codegreen},
    keywordstyle=\color{codeblue},
    numberstyle=\tiny\color{codegray},
    stringstyle=\color{codepurple},
    basicstyle=\ttfamily\scriptsize,
    breakatwhitespace=false,
    breaklines=true,
    captionpos=b,
    keepspaces=true,
    numbers=left,
    numbersep=5pt,
    showspaces=false,
    showstringspaces=false,
    showtabs=false,
    tabsize=2,
    frame=single,
    rulecolor=\color{coderule},
    extendedchars=true,
    literate={á}{{\'a}}1 {é}{{\'e}}1 {í}{{\'i}}1 {ó}{{\'o}}1 {ú}{{\'u}}1
             {Á}{{\'A}}1 {É}{{\'E}}1 {Í}{{\'I}}1 {Ó}{{\'O}}1 {Ú}{{\'U}}1
             {ã}{{\~a}}1 {õ}{{\~o}}1 {Ã}{{\~A}}1 {Õ}{{\~O}}1
             {ç}{{\c{c}}}1 {Ç}{{\c{C}}}1 {â}{{\^a}}1 {ê}{{\^e}}1
             {×}{{\ensuremath{\times}}}1 {—}{{\textemdash}}1
}

\lstset{style=pythonstyle}

% Ambientes Teoremas
\theoremstyle{definition}
\newtheorem{definicao}{Definição}[chapter]
\newtheorem{teorema}{Teorema}[chapter]
\newtheorem{lema}{Lema}[chapter]
\newtheorem{corolario}{Corolário}[chapter]
\newtheorem{proposicao}{Proposição}[chapter]
\newtheorem{exemplo}{Exemplo}[chapter]
\newtheorem{exercicio}{Exercício}[chapter]
\newtheorem{observacao}{Observação}[chapter]
\newtheorem{hipotese}{Hipótese}[chapter]

% Comandos Personalizados
\newcommand{\codigo}[1]{\texttt{#1}}
\usepackage{underscore}
\newcommand{\destaque}[1]{\textbf{#1}}
\newcommand{\opencode}{\textsc{OpenCode}}
\newcommand{\opencodeeco}{\textsc{OpenCode Ecosystem}}
\newcommand{\qualisA}{\textit{Qualis A1}}
\newcommand{\tdd}{\textsc{TDD}}
\newcommand{\sdd}{\textsc{SDD}}
\newcommand{\doilink}[1]{\href{https://doi.org/#1}{\texttt{DOI: #1}}}
\newcommand{\arxivlink}[1]{\href{https://arxiv.org/abs/#1}{\texttt{arXiv:#1}}}
\newcommand{\nivelzero}{$\bigstar$}
\newcommand{\nivelbasico}{$\bigstar\bigstar$}
\newcommand{\nivelintermediario}{$\bigstar\bigstar\bigstar$}
\newcommand{\nivelavancado}{$\bigstar\bigstar\bigstar\bigstar$}
\newcommand{\nivelphd}{$\bigstar\bigstar\bigstar\bigstar\bigstar$}

\newcommand{\doi}[1]{\doilink{#1}}
\newcommand{\arxiv}[1]{\arxivlink{#1}}

% Configurações ABNT
\setlength{\parindent}{1.5cm}
\setlength{\parskip}{0.2cm}
\SingleSpacing
